\documentclass[11pt]{amsart}
\usepackage{amsmath,amssymb,amsthm}
\usepackage[margin=1.1in]{geometry}
\usepackage[hidelinks]{hyperref}
\usepackage{array}
\hypersetup{pdftitle={Asymptotically optimal degree for the double Hall property},pdfauthor={John Tripodo}}
\usepackage[T1]{fontenc}
\usepackage{lmodern}

\newtheorem{theorem}{Theorem}
\newtheorem{lemma}[theorem]{Lemma}
\newtheorem{proposition}[theorem]{Proposition}
\newtheorem{corollary}[theorem]{Corollary}
\theoremstyle{definition}
\newtheorem{definition}[theorem]{Definition}
\theoremstyle{remark}
\newtheorem{remark}[theorem]{Remark}

\newcommand{\F}{\mathbb{F}}
\newcommand{\sn}{\operatorname{SN}}
\newcommand{\freg}{f_{\mathrm{reg}}}
\newcommand{\fbi}{f_{\mathrm{bi,reg}}}

\title{Asymptotically optimal degree for the double Hall property}
\author{John Tripodo}
\thanks{\copyright\ 2026 John Tripodo. Licensed under \href{https://creativecommons.org/licenses/by/4.0/}{CC BY 4.0}. ORCID: \href{https://orcid.org/0009-0006-9180-4953}{0009-0006-9180-4953}.}
\date{22 October 2026}

\begin{document}

\begin{abstract}
A bipartite graph with parts $X$ and $Y$ has the double Hall property if every
$A\subseteq X$ with $|A|\ge 2$ has at least $|A|$ vertices in $Y$ with two or more
neighbours in $A$. Chen, Lavrov, Ma, Su and Vandenbussche showed that such a graph with
$|X|=n$ has maximum degree at least $(1+o(1))\sqrt{2n}$, gave balanced regular
constructions along infinitely many orders of degree $n^{\log_{79}13+o(1)}\approx n^{0.587}$, and random balanced
constructions of degree $(1+o(1))\sqrt{2n\log n}$. We show that the counting bound is
asymptotically sharp: for every sufficiently large $n$ there is a balanced $D$-regular
bipartite graph with $n$ vertices in each part, with the double Hall property on
\emph{both} sides, and
\[
  D \le \sqrt{2n} + O\!\left(n^{21/80}\right).
\]
The construction is deterministic. It repairs the empty class rectangles of the classical
homology semibiplane by rook-type local graphs, then pads and completes the result to any
order. The error exponent is half of any admissible short-prime-interval exponent.
\end{abstract}

\maketitle

\section{Introduction}

Let $G$ be a bipartite graph with parts $X$ and $Y$. For $A\subseteq X$ write
\[
  \sn(A)=\{y\in Y:\ |N(y)\cap A|\ge 2\}
\]
for the \emph{super-neighbours} of $A$. Following Chen, Lavrov, Ma, Su and
Vandenbussche~\cite{CLMSV}, $G$ has the \emph{double Hall property} (on $X$) if
$|\sn(A)|\ge|A|$ for every $A\subseteq X$ with $|A|\ge 2$. The property is motivated by
cycle questions of Salia and others; see~\cite{CLMSV,BGJNP}.

For $n\ge2$, let $f(n)$ be the least maximum degree of a double-Hall graph with $|X|=n$, and let
$\freg(n)$ be the least $d$ for which a balanced $d$-regular graph with $n$ vertices in
each part has the property. Applying the definition to the pairs $\{x,x'\}$ shows that every
two vertices of $X$ have at least two common neighbours. Counting paths of length two from a
fixed $x$ gives
\begin{equation}\label{eq:lower}
  2(n-1)\le \Delta(\Delta-1),\qquad\text{so}\qquad f(n)\ge \tfrac12\bigl(1+\sqrt{8n-7}\bigr)
\end{equation}
\cite[Theorem~1.10]{CLMSV}. The same paper constructs balanced regular examples along infinitely many orders of degree
$n^{\log_{79}13+o(1)}$ \cite[Theorem~1.11]{CLMSV} and shows that random balanced graphs
of degree $(1+o(1))\sqrt{2n\log n}$ work \cite[Theorem~1.12]{CLMSV}.

Write $\fbi(n)$ for the least degree of a balanced regular graph with the double Hall
property on both $X$ and $Y$. Clearly $f(n)\le\freg(n)\le\fbi(n)$.

\begin{theorem}\label{thm:main}
As $n\to\infty$,
\[
  f(n),\ \freg(n),\ \fbi(n) \;=\; (1+o(1))\sqrt{2n}.
\]
More precisely, $\fbi(n)\le\sqrt{2n}+O(n^{21/80})$ for all sufficiently large $n$.
\end{theorem}

The exponent comes from prime gaps. We prove the following transfer statement.

\begin{theorem}\label{thm:transfer}
Let $\tfrac12\le\theta<1$. Suppose that for all sufficiently large real $x$ there is a prime
in $[x-x^\theta,x]$. Then $\fbi(n)\le\sqrt{2n}+O(n^{\theta/2})$, and a witnessing graph
can be computed deterministically in time polynomial in $n$.
\end{theorem}

Baker, Harman and Pintz~\cite{BHP} proved the hypothesis with $\theta=21/40$, which gives
Theorem~\ref{thm:main}. Their threshold is effectively determinable, although
we do not compute a numerical value here. Any improvement in short-interval prime results improves the error
term in the same way; for example, the preprint~\cite{Li} claims $\theta=0.52$.

Along the orders $N=(q^2-1)/2$ with $q$ an odd prime power, no prime-gap input is needed.

\begin{theorem}\label{thm:semibiplane}
Let $q\ge 1153$ be an odd prime power, $N=(q^2-1)/2$ and $g=(q-1)/2$. There is a balanced
regular bipartite graph with $N$ vertices in each part, with the double Hall property on both
sides, of degree
\[
  q+2\lfloor\sqrt g\rfloor+4 \;=\; \sqrt{2N}+O(N^{1/4}).
\]
\end{theorem}

The constant $1153$ is the first odd $q$ with $g=(q-1)/2\ge 576$ (Lemma~\ref{lem:rook-interp}).
For smaller $q$ one can use $K_{g,g}$ in the repair step, at the cost of a worse
second-order term.

\subsection*{Method}
The starting point is the classical homology semibiplane over $\F_q$ \cite{Moorhouse}. Its
signed incidence matrix has orthogonal rows, so by the strong-quadrangularity principle
of Severini~\cite{Severini,SeveriniSzollosi} every point set that is not contained in one
parallel class has enough super-neighbours (Section~\ref{sec:semibiplane}). The only
failures are sets inside a single parallel class, and each such class faces an empty block
class. We fill these $q+1$ empty rectangles with a balanced regular double-Hall graph of
degree $O(\sqrt q)$, built from rook graphs (Section~\ref{sec:rook}). This costs only a
lower-order term in the degree. To reach every order $n$ we duplicate whole parallel
classes, which raises each degree by an explicitly bounded amount, and then complete to a
regular graph using a sparse degree-completion lemma proved from Hall's theorem
(Section~\ref{sec:completion}).

\subsection*{What is not new}
Homology semibiplanes, strong quadrangularity, rook graphs, vertex twinning and degree
completion via flows or matchings are classical; we credit each where it is used. The
contribution is their combination for the double Hall degree problem, which closes the
gap between the counting bound~\eqref{eq:lower} and the best previous construction.

\section{Notation and two monotonicity facts}

All graphs are finite, simple and bipartite with a fixed bipartition $(X,Y)$. Adding edges
cannot destroy a super-neighbour, so the double Hall property is preserved under adding
edges. If $X$ is partitioned into classes, a set $A\subseteq X$ with $|A|\ge2$ either meets
two classes or lies in one class.

\begin{lemma}[class repair]\label{lem:repair}
Let $X$ be partitioned into classes. Let $B$ and $L$ be bipartite graphs on $(X,Y)$. Suppose
$|\sn_B(A)|\ge|A|$ whenever $|A|\ge2$ and $A$ meets two classes, and
$|\sn_L(A)|\ge|A|$ whenever $|A|\ge2$ and $A$ lies in one class. Then $B\cup L$ has the
double Hall property on $X$. If $B$ and $L$ are edge-disjoint, degrees in $B\cup L$ are
sums of degrees in $B$ and in $L$.
\end{lemma}

\section{Orthogonal supports}

\begin{lemma}[strong quadrangularity {\cite{Severini,SeveriniSzollosi}}]\label{lem:orth}
Let $W$ be a real matrix with rows indexed by $X$ and columns by $Y$. Let $A\subseteq X$,
and let $C(A)$ be the set of columns with at least two nonzero entries in the rows of $A$.
Suppose the rows of $A$ are pairwise orthogonal, and every row of $A$ shares a nonzero
column with another row of $A$. Then $|C(A)|\ge|A|$.
\end{lemma}

\begin{proof}
Restrict the rows of $A$ to the columns in $C(A)$. A column outside $C(A)$ has at most one
nonzero entry in the rows of $A$, so deleting it changes no inner product between distinct
rows of $A$. Hence the restricted rows are still pairwise orthogonal. Every row of $A$ shares a
nonzero column $y$ with another row of $A$; then $y\in C(A)$, so the restricted row is
nonzero. Pairwise orthogonal nonzero vectors are linearly independent, so
$|A|\le\dim\mathbb{R}^{C(A)}=|C(A)|$.
\end{proof}

\section{The homology semibiplane}\label{sec:semibiplane}

Let $q$ be an odd prime power and $\F=\F_q$. Let $P$ be the set of nonzero vectors of
$\F^2$ modulo $u\sim-u$, so $|P|=(q^2-1)/2=:N$. Take $X$ and $Y$ to be two disjoint copies
of $P$ and join $[u]\in X$ to $[v]\in Y$ when $u\cdot v\in\{1,-1\}$. This relation does not
depend on the chosen representatives. This is the Desarguesian homology semibiplane
\cite[\S\S1--2]{Moorhouse}.

\begin{lemma}\label{lem:sb-basic}
\begin{enumerate}
\item The graph is $q$-regular.
\item Partition $X$ into the $q+1$ parallel classes $C_\ell=\{[u]:u\in\ell\}$, one for
each one-dimensional subspace $\ell$; each has $g=(q-1)/2$ elements. Partition $Y$ into
$E_\ell=\{[v]:v\in\ell^\perp\}$. There are no edges between $C_\ell$ and $E_\ell$.
\item Every $y\in Y$ has at most one neighbour in each class $C_\ell$, and every $x\in X$
has at most one neighbour in each class $E_\ell$.
\end{enumerate}
\end{lemma}

\begin{proof}
(1) For $u\ne0$, each affine line $\{v:u\cdot v=s\}$ has $q$ points, and the lines for
$s=1$ and $s=-1$ are disjoint and swapped by $v\mapsto-v$. So exactly $q$ classes $[v]$
are adjacent to $[u]$. The same holds with the roles of $u$ and $v$ exchanged.
(2) A line through the origin has $q-1$ nonzero vectors, forming $(q-1)/2$ classes. If
$u\in\ell$ and $v\in\ell^\perp$ then $u\cdot v=0$.
(3) If $u$ and $\lambda u$ both satisfy $u\cdot v,\ \lambda u\cdot v\in\{\pm1\}$, then
$\lambda=\pm1$, so $[\lambda u]=[u]$. The second statement is symmetric.
\end{proof}

\begin{proposition}\label{prop:sb-hall}
If $A\subseteq X$ meets two parallel classes, then $|\sn(A)|\ge|A|$. The same holds on $Y$.
\end{proposition}

\begin{proof}
Choose representatives and define $W_{[u],[v]}=1$ if $u\cdot v=1$, $W_{[u],[v]}=-1$ if
$u\cdot v=-1$, and $0$ otherwise. Its nonzero pattern is the incidence. Write
$\sigma(t)\in\{1,-1,0\}$ for this sign of $t\in\F$, so $\sigma(-t)=-\sigma(t)$.

\emph{Orthogonality.} Let $[a]\ne[b]$. The product $\sigma(a\cdot v)\sigma(b\cdot v)$ is
even in $v$, so the inner product of rows $[a]$ and $[b]$ is half the sum of this product
over all $v\in\F^2$. If $a,b$ are linearly independent, $v\mapsto(a\cdot v,b\cdot v)$ is a
bijection of $\F^2$. The sum then equals $(\sum_{s\in\F}\sigma(s))^2=0$. If $b=\lambda a$,
both factors are nonzero only when $\lambda=\pm1$, which would give $[a]=[b]$. So every
term vanishes.

\emph{Sharing.} If $A$ meets two classes, every $[a]\in A$ has some $[b]\in A$ in another
class. Then $a,b$ are independent, so some $v$ has $a\cdot v=b\cdot v=1$, and column $[v]$
is nonzero in both rows.

Lemma~\ref{lem:orth} now gives $|\sn(A)|=|C(A)|\ge|A|$. The incidence is symmetric under
exchanging $X$ and $Y$, and the classes $E_\ell$ correspond to the classes
$C_{\ell^\perp}$, so the same argument applies on $Y$.
\end{proof}

\section{Rook graphs and matched duplication}\label{sec:rook}

\begin{lemma}[rook cross]\label{lem:rook}
For $m\ge2$ let $X=Y=[m]^2$ (disjoint copies), with $(a,b)\in X$ adjacent to $(i,j)\in Y$
when $a=i$ or $b=j$. This graph is $(2m-1)$-regular, contains the diagonal matching
$(a,b)\leftrightarrow(a,b)$, and has the double Hall property on both sides.
\end{lemma}

\begin{proof}
Regard $A\subseteq X$ as a set of occupied cells. An occupied cell is itself a
super-neighbour unless it is alone in both its row and its column; call such cells
\emph{isolated}. All other occupied cells are distinct super-neighbours. If there are $k\ge2$
isolated cells $(a_1,b_1),\dots,(a_k,b_k)$, their rows and columns are distinct, and the
unoccupied cells $(a_i,b_{i+1})$ (indices mod $k$) are $k$ further distinct
super-neighbours. If there is exactly one isolated cell $(a,b)$, take another occupied cell
$(c,d)$. The cell $(a,d)$ is unoccupied, since $(a,b)$ is alone in its row, and it sees both.
So it is one more super-neighbour. In all cases $|\sn(A)|\ge|A|$. The graph is symmetric.
\end{proof}

\begin{lemma}[matched duplication]\label{lem:dup}
Let $H$ be balanced, with $X=\{x_i\}$, $Y=\{y_i\}$ and every $x_iy_i$ an edge, and with the
double Hall property on both sides. For a set $T$ of indices, duplicate $x_i$ and $y_i$ for
each $i\in T$, copying all incidences, including those between duplicates. The result has
the double Hall property on both sides. If $H$ is $d$-regular, the new degree of a copy of
$x_i$ is $d+|N_H(x_i)\cap\{y_j:j\in T\}|$, and symmetrically on $Y$.
\end{lemma}

\begin{proof}
Let $A$ be a set of new left vertices, $B$ its set of original indices, and $k$ the number
of doubled indices with both copies in $A$, so $|A|=|B|+k$. If $|B|\ge2$, every
super-neighbour of $B$ in $H$ is a super-neighbour of $A$. Each doubled index $i$ also gives
the duplicate of $y_i$, which sees both copies of $x_i$; these are $k$ further distinct
super-neighbours. If $|B|=1$, then $A$ is one doubled pair, which has two common neighbours
$y_i,y_i'$. The other side is symmetric.
\end{proof}

\section{Sparse regular completion}\label{sec:completion}

\begin{lemma}[completion]\label{lem:completion}
Let $H$ be a balanced bipartite graph with $n$ vertices per side and maximum degree at most
$D\le n$. Let $\alpha_x=D-\deg_H x$ and $\beta_y=D-\deg_H y$, with total deficit
$T_0=\sum\alpha_x=\sum\beta_y$, largest deficits $a=\max\alpha$, $b=\max\beta$, and largest
degrees $r=\max_X\deg_H$, $c=\max_Y\deg_H$. If
\[
  T_0\ \ge\ ac+br+2ab,
\]
then $H$ is a subgraph of a simple $D$-regular bipartite graph on the same parts.
\end{lemma}

\begin{proof}
We need a subgraph $F$ of the bipartite complement $\bar H$ with $\deg_F x=\alpha_x$ and
$\deg_F y=\beta_y$.

\emph{Reduction to a perfect matching.} Build a bipartite graph $\Gamma$ (Tutte's gadget).
Its left part is $\alpha_x$ copies of each $x$, together with one vertex $e^Y$ for each edge
$e$ of $\bar H$. Its right part is $\beta_y$ copies of each $y$, together with one vertex
$e^X$ for each $e$. Join every copy of $x$ to $e^X$ for each $e\ni x$. Join $e^Y$ to $e^X$.
Join $e^Y$ to every copy of $y$ for each $e\ni y$. Both parts have $T_0+|E(\bar H)|$
vertices. In a perfect matching, put $e=xy$ into $F$ exactly when $e^X$ is matched to a
copy of $x$. This happens exactly when $e^Y$ is matched to a copy of $y$. Then $F$ has the
required degrees.

\emph{Hall's condition.} Let $S$ be a set of left vertices. Let $U$ be the set of $x$ having a
copy in $S$, and let $S_2$ be the set of gadget edges $e^Y$ in $S$, with $Y$-endpoints $V$.
Then $N(S)$ contains every $e^X$ with $e$ at $U$ or $e\in S_2$, and every copy of every
$y\in V$. Enlarging $S_2$ to all edges of $\bar H$ at $V$ does not change $V$, so in
checking $|N(S)|\ge|S|$ we may assume this. With $V'=Y\setminus V$ the requirement becomes
\begin{equation}\label{eq:cut}
  \alpha(U)+\beta(V')-T_0\ \le\ e_{\bar H}(U,V')\qquad(U\subseteq X,\ V'\subseteq Y).
\end{equation}
If $|V'|\ge r+a$ then $e_{\bar H}(U,V')\ge|U|(|V'|-r)\ge a|U|\ge\alpha(U)$, which
dominates the left side. If $|U|\ge c+b$ the symmetric bound applies. Otherwise
$\alpha(U)+\beta(V')\le a(c+b)+b(r+a)\le T_0$. So~\eqref{eq:cut} holds, and Hall's theorem
gives the perfect matching.
\end{proof}

\begin{remark}
The inequality \eqref{eq:cut} is the cut condition of the integral flow formulation, which
also gives a polynomial-time algorithm. The reduction through Hall's theorem is used here
because it is elementary and is the route taken in the formalization.
\end{remark}

\begin{lemma}[rook interpolation]\label{lem:rook-interp}
For every $g\ge576$ there is a balanced regular bipartite graph on $g+g$ vertices with the
double Hall property on both sides, containing a perfect matching, of degree
$2\lfloor\sqrt g\rfloor+4$.
\end{lemma}

\begin{proof}
Let $m=\lfloor\sqrt g\rfloor$ and $t=g-m^2\le 2m$. In the rook-cross graph on $[m]^2$
duplicate $t$ distinct cells from the two cyclic diagonals $\{(i,i)\}\cup\{(i,i+1)\}$.
Every row and every column of $[m]^2$ then contains at most two selected cells, so each
rook neighbourhood contains at most four selected cells. An original vertex
and its duplicate have the same original neighbourhood, so each gains at most
four edges from the duplication. By Lemma~\ref{lem:dup}, all degrees lie in
$[2m-1,2m+3]$. Complete to $D=2m+4$ by Lemma~\ref{lem:completion}. Here $T_0\ge g$,
$a,b\le5$ and $r,c\le 2m+3$, so the condition holds once $m^2\ge 20m+80$, that is, $m\ge24$.
Completion adds edges, so both double Hall properties and the diagonal
matching survive by monotonicity.
\end{proof}

\section{Proofs of the main theorems}

\begin{proof}[Proof of Theorem~\ref{thm:semibiplane}]
Take the semibiplane of Section~\ref{sec:semibiplane}. For each $\ell$ identify $C_\ell$ and
$E_\ell$ with the two parts of the graph of Lemma~\ref{lem:rook-interp} on $g$ vertices,
of degree $h=2\lfloor\sqrt g\rfloor+4$, and insert it into the empty rectangle
$C_\ell\times E_\ell$. The inserted graphs are edge-disjoint from the base and from each
other. So the result is $(q+h)$-regular. Sets meeting two classes are handled by
Proposition~\ref{prop:sb-hall}, sets inside one class by the local graph, and
Lemma~\ref{lem:repair} applies on both sides.
\end{proof}

\begin{proof}[Proof of Theorem~\ref{thm:transfer}]
Let $Q=\sqrt{2n+1}$. By the hypothesis there is an odd prime $q\le Q$ with
$Q-q=O(Q^\theta)$. Put $N=(q^2-1)/2\le n$, $g=(q-1)/2$ and $t=n-N=(Q-q)(Q+q)/2=O(q^{1+\theta})$.
Let $G$ be the $d$-regular graph of Theorem~\ref{thm:semibiplane}, $d=q+O(\sqrt q)$. The
local diagonals give a perfect matching $x\leftrightarrow\pi(x)$ with $\pi(C_\ell)=E_\ell$.

\emph{Class-block padding.} Write $t=kg+s$ with $0\le s<g$, and order the parallel classes
$\ell_1,\ell_2,\dots$. Let $T$ consist of all indices in $C_{\ell_1},\dots,C_{\ell_k}$
together with $s$ indices of $C_{\ell_{k+1}}$. This needs $k+1\le q+1$, which holds for
large $n$ since $t=o(N)$. Duplicate $T$ by Lemma~\ref{lem:dup}. Consider any original
$x\in C_\ell$. By Lemma~\ref{lem:sb-basic}(3) its base neighbours meet each block class in at
most one vertex. Its local neighbours lie in $E_\ell$. The duplicated columns lie in at most
$k+1$ block classes. Hence the load of $x$ is at most
\[
  R := k+1+h,
\]
and likewise on $Y$. All degrees now lie in $[d,d+R]$, and both sides have the double Hall
property.

\emph{Completion.} Set $D=d+R+1$. Then $T_0\ge n$, $a,b\le R+1$ and $r,c\le d+R$, so
Lemma~\ref{lem:completion} applies once $n\ge 2(R+1)(d+R)+2(R+1)^2$. Since
$k\le t/g=O(q^\theta)$ and $h=O(\sqrt q)$, we have $R=O(q^\theta)$ and
$(R+1)(d+R)=O(q^{1+\theta})=o(n)$, so the condition holds for large $n$. Completion only adds edges, and therefore
preserves the double Hall property on both sides. The final degree is
\[
  D=q+O(q^\theta)=\sqrt{2n}+O(n^{\theta/2}).
\]
Every step is explicit: finding $q$, finite-field arithmetic, the rook graphs, duplication,
and two completions computed by integral max-flow.
\end{proof}

\begin{proof}[Proof of Theorem~\ref{thm:main}]
Combine Theorem~\ref{thm:transfer} with $\theta=21/40$ \cite{BHP} and the lower
bound~\eqref{eq:lower}, which applies already to one side.
\end{proof}

\section{Remarks and open problems}

\begin{enumerate}
\item The order of magnitude of the repair degree is optimal for this base. If edges are added to the
homology semibiplane so that every vertex gains at most $r$ edges and double Hall holds on
both sides, then $r(r+1)\ge q-3$. Indeed, same-class pairs have no common neighbour in the
base, and each pair needs two. At each column, new--old common-neighbour witnesses
contribute at most $r$, and new--new witnesses at most $\binom r2$. Thus
$N(g-1)\le N(r+\binom r2)$, giving the asserted bound. So the $N^{1/4}$-order second term in
Theorem~\ref{thm:semibiplane} cannot be improved by local repair of this base.
\item Exact values of $\freg(n)$ remain open. Equality in \eqref{eq:lower} forces a
biplane-type structure. Infinite families of biplanes are not known.
\item Our graphs say nothing directly about the cycle conjectures that motivate the
double Hall property \cite{CLMSV,BGJNP}.
\end{enumerate}

\section*{Formal verification}
The counting, construction and existence statements mapped below have a Lean~4
formalization with Mathlib (commit \texttt{0ed37c9258f1}). The prime-interval input
\cite{BHP} is an explicit hypothesis of the formal transfer theorem, not a result proved
inside the companion. The algorithmic complexity remarks and the augmentation lower bound
in the first item of the preceding section are written arguments, not claims of formal
coverage by this companion. Every declaration in the recorded endpoint audit depends only
on \texttt{propext}, \texttt{Classical.choice} and \texttt{Quot.sound}.

\medskip
\begin{tabular}{>{\raggedright\arraybackslash}p{0.43\textwidth}>{\raggedright\arraybackslash}p{0.5\textwidth}}
Paper & Lean declaration (namespace \texttt{DoubleHall})\\\hline
Lower bound \eqref{eq:lower} & \texttt{pair\_count\_lower\_bound}, \texttt{hasRegularBoth\_lower}\\
Lemma~\ref{lem:repair} & \texttt{doubleHallLeft\_union\_of\_classes}, \texttt{localRel\_hall}\\
Lemma~\ref{lem:orth} & \texttt{card\_le\_card\_multiCols}\\
Lemma~\ref{lem:sb-basic}, Prop.~\ref{prop:sb-hall} & \texttt{Semibiplane.degL\_inc}, \texttt{two\_mul\_card\_cls},
  \texttt{not\_inc\_of\_cls\_eq\_clsY}, \texttt{inc\_cls\_unique}, \texttt{W\_orthogonal},
  \texttt{card\_le\_superNbrs\_of\_not\_inOneClass}\\
Lemma~\ref{lem:rook} & \texttt{Rook.doubleHallLeft}, \texttt{Rook.degL\_eq}\\
Lemma~\ref{lem:dup} & \texttt{Dup.doubleHallLeft}, \texttt{Dup.degL\_dup}\\
Lemma~\ref{lem:completion} & \texttt{exists\_factor}, \texttt{cutCondition\_of\_sparse},
  \texttt{exists\_regular\_supergraph}\\
Lemma~\ref{lem:rook-interp} & \texttt{RookInterp.exists\_local}\\
Theorem~\ref{thm:semibiplane} & \texttt{semibiplane\_order\_theorem}\\
Class-block padding & \texttt{Semibiplane.padding\_theorem}\\
Theorem~\ref{thm:transfer} & \texttt{transfer\_theorem} (explicit constant $33$)\\
\end{tabular}

\medskip
The formal statement of Theorem~\ref{thm:transfer} is: if $\tfrac12\le\theta<1$ and every
sufficiently large real $x$ has a prime in $[x-x^\theta,x]$, then for all large $n$ there is a
balanced $D$-regular graph with $n$ vertices per side, double Hall on both sides, and
$D\le\sqrt{2n}+33\,n^{\theta/2}$. The algorithmic remarks in the proofs are not formalized: the
formal theorems are existence statements.

\section*{Use of AI tools}
AI tools were used for exploration, drafting, computation and formalization. The author
is responsible for the mathematical statements, proofs, source comparisons and final text.

\begin{thebibliography}{99}
\raggedright
\bibitem{BHP} R.~C.~Baker, G.~Harman and J.~Pintz, The difference between consecutive
primes, II, \emph{Proc. London Math. Soc.} (3) 83 (2001), 532--562.
\bibitem{BGJNP} J.~Bar\'at, A.~Grzesik, A.~Jung, Z.~L.~Nagy and D.~P\'alv\"olgyi, The double
Hall property and cycle covers in bipartite graphs, \emph{Discrete Math.} 347 (2024), 114079.
\bibitem{CLMSV} G.~Chen, M.~Lavrov, Y.~Ma, Y.~Su and J.~Vandenbussche, Bipartite graphs with
the double Hall property, \emph{Frontiers in Combinatorics and Number Theory} 1 (2026), 38--59,
doi:10.3934/fcnt.2026005.
\bibitem{Li} R.~Li, The number of primes in short intervals and numerical calculations for
Harman's sieve, \mbox{arXiv:2308.04458v8} (2025), preprint.
\bibitem{Moorhouse} G.~E.~Moorhouse, On the construction of finite projective planes from
homology semibiplanes, \emph{European J. Combin.} 11 (1990), 589--600.
\bibitem{Severini} S.~Severini, On the digraph of a unitary matrix, \emph{SIAM J. Matrix Anal.
Appl.} 25 (2003), 295--300; \mbox{arXiv:math/0205187}.
\bibitem{SeveriniSzollosi} S.~Severini and F.~Sz\"oll\H{o}si, A further look into
combinatorial orthogonality, \emph{Electron. J. Linear Algebra} 17 (2008), 376--388.
\end{thebibliography}

\end{document}
